\documentclass[12pt, a4paper]{article}
\usepackage[english]{babel}
\usepackage{graphicx}
\usepackage{url}
\renewcommand{\title}{On the use of knowledge-based media to enhance the social life of elderly}
\newcommand{\authorA}{Harrie Oosterhuis}
\newcommand{\studnummerA}{10196129}
\newcommand{\authorB}{Dasyel Willems}
\newcommand{\studnummerB}{10172548}
\newcommand{\authorC}{Carla Groenland}
\newcommand{\studnummerC}{10208429}
\newcommand{\authorD}{Fabian Voorter}
\newcommand{\studnummerD}{10218807}
\newcommand{\authorE}{Jelle van Assema}
\newcommand{\studnummerE}{10200363}
\newcommand{\einddatum}{\today}
\newcommand{\content}{Report Kennisgebaseerde media}

\usepackage[pdftex,pagebackref,colorlinks,linkcolor=black,urlcolor=blue,citecolor=black]{hyperref}
\begin{document}

\begin{titlepage}
\vspace*{10pt}
\begin{center}
{\LARGE \textbf{\bf \title}\par}\bigskip
{\Large \authorA, \authorB, \authorC, \authorD, \authorE
\par}\bigskip
{\Large \einddatum\par}\bigskip\bigskip
\vspace{\stretch{1}}
\vspace{\stretch{1}}
Bachelor Kunstmatige Intelligentie \par\smallskip
Faculteit der Natuurwetenschappen, Wiskunde en Informatica\par\smallskip
Universiteit van Amsterdam\par\bigskip
\includegraphics[width=0.075\hsize]{uvalogo}
\end{center}
\end{titlepage}

\begin{titlepage}
\noindent
\begin{minipage}{0.85\hsize}
\section*{Abstract}
This report describes the possibilities of using knowledge-based media systems to enhance the social life of elderly. 
\end{minipage}

\vfill

\noindent Title: \title\\
Authors: \authorA, \studnummerA\\
\hphantom{Authors: }\authorB, , \studnummerB\\
\hphantom{Authors:\ }\authorC,  \studnummerC\\
\hphantom{Authors:\ }\authorD,  \studnummerD\\
\hphantom{Authors:\ }\authorE, \studnummerE\\
Date: \einddatum

\bigskip\noindent
Faculteit der Natuurwetenschappen, Wiskunde en Informatica\\
Universiteit van Amsterdam\\
Science Park 904, 1098 XH Amsterdam\\
\url{http://www.science.uva.nl/home.cfm}
\end{titlepage}

\tableofcontents 
\newpage
\section{Introduction}
A complaint amongst older people is that they are bored or lonely; they miss the social interaction younger people automatically have from sport, work or school. In this report, we discuss a system which tries to enhance the life of elderly by providing them with entertainment and by helping them with their social life. The system is aimed to respect privacy and the needed rest of the elderly and is equipped with a user-friendly interface.

The system reduces the boredom en loneliness of the elderly by providing them with user-based entertainment. [EXPLANATION OF MEDIA]

Furthermore, the system enhances the social life of the elderly by planning several social activities for the elderly based on their personal needs. [TYPES OF ACTIVITIES?]

The communication of the elderly is through [EXPLANATION INTERFACE]

The system adapts to every user by gathering a lot of data about the person, both about his life history, interests and family and about his interactions with the system. As an example, the system tries to recognize patterns in which type of entertainment a person in most pleased with and on which occasions. Through large amounts of data, the system gets to know each user and behaves appropriately.

\section{Goal and system tasks}

\section{Literature review}

\section{System description}

\section{Discussion of approach}

\section{Conclusion}
\begin{thebibliography}{x}
\addcontentsline{toc}{section}{Bibliografy} 

\end{thebibliography}

\end{document}